% Source-derived chapter 02: Conventions
% Original source: standard-conjectures-abelian-fourfolds
\chapter{Conventions}
\label{standard-conjectures-abelian-fourfolds-chapter-02}
\source{standard-conjectures-abelian-fourfolds}

\begin{remark}[Source section]
\label{standard-conjectures-abelian-fourfolds-chapter-02-source}
\source{standard-conjectures-abelian-fourfolds}
\source{standard-conjectures-abelian-fourfolds}
This chapter reproduces the corresponding section of the retrieved
LaTeX source, including its definitions, statements, examples, and
proofs. It has not yet been translated into Lean.
\notready
\end{remark}

% The source section title was: Conventions
Throughout the paper we will use the following conventions.
\begin{enumerate}
\item  A \textbf{variety} will mean a smooth, projective and geometrically connected scheme over a base  field  $k$. 



\item Let  $k$ be a base field and $F$ be a field of coefficients of characteristic zero, we will denote by \[\CHM(k)_{F}\] the category of  \textbf{Chow motives} over $k$ with coefficients in $F$. (The subscript $F$ may sometimes be omitted.) For generalities, we refer to \cite{Andmot} in particular to \cite[Definition 3.3.1.1]{Andmot}  for the notion of Weil cohomology and to \cite[Proposition 4.2.5.1]{Andmot}  for the associated realization functor. 

We will denote by
\[\mathfrak{h} : \Var_k\longrightarrow \CHM(k)_{F}^{\op}\]
the functor associating to each variety its motive.

We will work also with \textbf{homological and numerical motives}. The motive of a variety $X$ in these categories will still be denoted by $\mathfrak{h} (X)$. We will denote by 
\[\NUM(k)_F\]
the category of numerical motives over $k$ with coefficients in $F$. 

Finally, we will need also to work with  \textbf{relative motives} as defined for example in \cite[\S 5.1]{OS2}.  The relative situation that will appear in this paper will always be over the ring of integers $W$ of a $p$-adic field. We will use analogous notation as in the absolute setting, for example Chow motives over $W$ will be denoted by \[\CHM(W).\]

\item By \textbf{classical realizations}\footnote{We tried to distinguish the properties which hold true for any realization and those which are known only for classical realizations.  In any case, the main results of the paper only make use of classical realizations and the reader can safely think only about them.} we will mean Betti realization
\[R_B : \CHM(\C)_{\Q}\longrightarrow \GrVec_\Q,\]
  $\ell$-adic realization (when $\ell$ is invertible in $k$)
\[R_\ell : \CHM(k)_{\Q}\longrightarrow \GrVec_{\Q_\ell}\]
 de Rham realization (when $k$ is of characteristic zero)
\[R_\dR : \CHM(k)_{\Q}\longrightarrow \GrVec_{k}\]
and crystalline realization (when $k$ is perfect and of positive characteristic)
\[R_\crys : \CHM(k)_{\Q}\longrightarrow \GrVec_{\Frac (W(k))}.\]
 For generalities see  \cite[\S 3.4.2 and Proposition 4.2.5.1]{Andmot}.


The realization functor are sometimes used in their enriched form, i.e. $R_B$ provides Hodge structures, $R_\ell$ provides Galois representations, $R_\dR$ provides filtered vector spaces and $R_\crys$ provides absolute Frobenius actions, see \cite[\S 7.1]{Andmot}.

\item\label{Hyodo-Kato} Let $W$ be the ring of integers  of a $p$-adic field $K$ with residue field $k$ (then $\Frac (W(k))$ is the maximal unramified subfield of $K$). We will denote by  \[\MOD\] the category  of admissible filtered $\varphi$-modules \cite{BertOgus,HyodoKato,NiziFred}.
\source{standard-conjectures-abelian-fourfolds}

Recall that an element in $\MOD$ is, in particular, a finite dimensional  $\Frac (W(k))$-vector space $V$, together with a decreasing filtration $\Fil^*$ on $V\otimes_{\Frac (W(k))} K$ and a $(\Frac (W(k))/\Q_p)$-semilinear endomorphism $\varphi$ of $V$, usually called absolute Frobenius, verifying some compatibilities between $\Fil^*$ and $\varphi$. 

There is a \textbf{Hyodo-Kato} realization functor \[R_{\HK}: \CHM(W)\rightarrow \MOD\]
such that, after forgetting the absolute Frobenius, it coincides with de Rham realization of the generic fiber $R_{\HK}(\cdot)\otimes K=R_\dR(\cdot_{\vert K})$  and, after forgetting the filtration, it coincides with  crystalline realization of the special fiber $R_{\HK}(\cdot) = R_{\crys}(\cdot_{\vert k})$. 
\item The \textbf{unit object} in one of the above categories of motives (Chow, homological or numerical motives) over a base $S$ is the motive $\mathfrak{h}(S)$ and it will be denoted by
\[\one \coloneqq \mathfrak{h}(S).\]
When $S=\Spec(k)$ we have the identification
\[\End(\one)=\Q.\]
\item An \textbf{algebraic class} of a motive $T$ is a map \[\alpha\in\Hom(\one, T).\]
The realization of an algebraic class gives a map
\[R(\alpha) : F \longrightarrow R(T),\]
where $F$ is the field of coefficients of the realization. 
The map $R(\alpha)$ is characterized by its value at $1\in F$. An element of the cohomology group $R(T)$ of this form
\[R(\alpha)(1) \in R(T)\]
will be called an algebraic class of  $R(T)$. By abuse of language, an algebraic class of the cohomology group $R(T)$ may also be called an algebraic class of the motive $T$.


Finally, we  may sometimes ignore Tate twists and say that $T$ has algebraic classes when $T(n)$ does. For example, for a variety $X$, the cohomology group $H^{2n}(X)$ may have algebraic classes, although, strictly speaking, these classes belong to $H^{2n}(X)(n).$ Hopefully this will not bring confusion as we will work with motives whose realization is concentrated in one single cohomological degree.
 

\item An algebraic class $\alpha$ of a motive $T$ is \textbf{numerically trivial} if, for all $\beta\in\Hom(T,\one)$, the composition
$\beta \circ \alpha $  is zero. Two algebraic classes are numerically equivalent if their difference is numerically trivial. (This recovers the classical definition \cite[7.1.2]{AK}.)

\item Let $X$ be a variety. The space of algebraic cycles on $X$ with coefficients in $F$ modulo numerical equivalence will be denoted by \[\mathcal{Z}^{n}_{\num}(X)_{F}= \Hom_{\NUM(k)_F}(\one, \mathfrak{h}(X)(n)).\]
Similarly, for a fixed Weil cohomology, $\mathcal{Z}^{n}_{\hom}(X)_{F}$ will denote the space of algebraic cycles on $X$ with coefficients in $F$ modulo homological equivalence.



\end{enumerate}
